% R. Brown's translations of two letters to L. Breen, with his note and comments (cote 134-2): archive cote 134-2, site page 136 (PDF page 137).
% translation of the letter of 17-19 July 1975. Draft compared: ps1.tex, lines 1570-1585.
% AI-transcribed by Claude Opus 5.5 (claude-opus-5-5, adaptive thinking, effort high), Message Batches API, batch brown_full (msgbatch_01Md3i5yK1PxkpJn5GHfEExE), 2026-11-13; not yet checked by eye.
\scanpage{134-2/136}
\mspage{- 9 -}

that we must make more explicit the nature of the left exact
functor $\phi$ . I think that this should be possible, in the context
of \uline{schemes} in which I am placed (technically rather sophisticated),
on introducing an ``étale tubular neighbourhood'' of $X_0$ in $X_1$
(which is a very satisfying topos, but not associated to a sheaf).
But this technical construction is only a paraphrase of a topol-%
ogical intuition extraordinarily simple, which one can make
explicit, on supposing to fix the ideas that the base field is
$\mathbb{C}$ and so one profits by working with locally compact spaces.
The topological idea behind the hypothesis of equi-singularity is
that there exist a \uline{tubular} \uline{neighbourhood} $T$ of $X_0$ in $X$
retracting onto $X_0$ and such that the pair $(X_0, T)$ over $X_0$
should be a locally trivial bundle, i.e. that $T\backslash X_0$ is locally
trivial over $X_0$ . In fact if $\partial T$ is the ``boundary'' of $T$ ,
which also should be a locally trivial bundle on $X_0$ , then $T/X$
is the bundle with fibres
$(\simeq (\partial T \times I) \amalg_{\partial T} X_0$ , where $I = [0, 1]$ , $\partial T \to \partial T \times I$ . is defined
by $x \longmapsto (x, 1)$ , and $\partial T \to X_0$ is the projection) then
$\mathring{T} = T\backslash X_0 \simeq \partial T \times [0, 1]$ is $X_0$-homotopic to $\partial T$ . If $X_0$ and
$X_1$ are non singular, then the same is true for $\mathring{T}$ and $\partial T$ ,
which are then topologically smooth fibrations on $X_0$ . Moreover,
putting $\tilde{X}_1 = X_1\backslash\mathring{T}$ , the inclusion $\tilde{X}_1 \to X_1$ is a homotopy equivalence
and $X$ can be recovered, up to homeomorphism, from the diagram
of spaces
\diagram{\begin{tikzcd}
(\mathring{T} = T\backslash X_0 \cong)\partial T \ar[r, "j", "\text{inclusion}"'] \ar[d, "\text{fibration}"', "p"] & \tilde{X}_1(\cong X_1) \\
X_0 &
\end{tikzcd}}

\margin{137}
